% !TEX root = ../../../tree_notes.tex
% =====================================================================
%  15.3  Prefix/suffix merge over children   sec:c20-prefixsuffix
%  Exports: def:c20-pre-suf  thm:c20-pre-suf-linear  lst:c20-presuf
%           fig:c20-presuf  ex:c20-presuf-run  prob:c20-presuf-solved
%           lst:c20-960e-soln
%  The obstruction 15.1/15.2 dodged by invertibility: when Phi is max / min /
%  signed-product / or, you cannot get "all children except c" by subtracting.
%  Prefix/suffix precompute over v's children buys that aggregate in O(1) per
%  child after O(deg) work; total O(sum deg) = O(n). Exclusive segment tree is
%  the forward-cited alternative.
%  Anchor: CF 960E Alternating Tree -- sum over ALL directed paths of the
%  alternating-sign node-value sum, mod 1e9+7. Reroot with a per-edge SIGN
%  FLIP; the per-node score multiplies two branch aggregates, so the general
%  "everything-but-this-branch" combine is done by prefix/suffix.
%  Overlay hand-computed on T; official n=4 sample gives 40.
% =====================================================================
\section{Prefix/Suffix Merge: When the Combine Has No Inverse}\label{sec:c20-prefixsuffix}

\begin{previouslybox}
\begin{enumerate}[nosep, leftmargin=1.6em]
  \item \secref{c20-twodfs} and \secref{c20-updown} recovered a child's
        siblings by $\mathit{down}[v]\ominus\mathit{down}[c]$. What property of
        $\Phi$ made that legal? (invertibility -- a sum or an XOR)
  \item For $\Phi=\max$, name the value of $\max(5,3,\mathbf{7},2)$ with the
        $\mathbf{7}$ removed. Can you get it by ``$\max(\cdot)$ minus
        $7$''? (the answer is $5$; no, subtraction is meaningless for $\max$)
  \item Ch.~13's in/out independent-set pair (\thmref{c10-mis-tree}) uses
        $\max$. If you had to report that pair for \emph{every} root, which
        step of the up pass would break without an inverse? (forming
        "all children but $c$" at the parent)
  \item How many children does a node of degree $d$ have in the rooted view,
        and what is $\sum_v \deg(v)$? ($d-1$ or $d$; the sum is $2(n-1)$)
\end{enumerate}
\end{previouslybox}

\begin{topicintro}
Sections 15.1 and 15.2 were lucky: their combines were sums, so ``the fold over
all children except $c$'' was one subtraction. Take that luck away -- let the
combine be a $\max$, a $\min$, a product with a sign, or a bitwise $\mathit{or}$
-- and the subtraction trick dies, because those operators have no inverse.
Recomputing the sibling fold per child costs $\bigO{\deg(v)}$ each, i.e.\
$\bigO{\deg(v)^2}$ per node and $\bigO{n^2}$ on a star. The repair is a two-line
idea: precompute a \term{prefix} fold and a \term{suffix} fold of the children
list, then ``everything except child $i$'' is one combination of a prefix and a
suffix. It costs $\bigO{\deg(v)}$ per node, $\bigO{n}$ overall
(\thmref{c20-pre-suf-linear}). This section builds that device and spends it on
CF~960E: \emph{the sum, over all directed paths, of the path's alternating-sign
node-value sum}, where the reroot carries a per-edge sign flip.
\end{topicintro}

\subsection{The Prefix/Suffix Aggregate}

\begin{definition}{Exclusive prefix/suffix combine over a node's children}{c20-pre-suf}
Let $v$ have children $c_0, c_1, \dots, c_{k-1}$ (in any fixed order) with child
values $x_i = \mathit{down}[c_i]$. Build two arrays of length $k{+}1$:
\[
  \mathit{pre}[0] = \mathcal{I}, \qquad \mathit{pre}[i{+}1] = \Phi\bigl(\mathit{pre}[i],\, x_i\bigr),
  \qquad \mathit{suf}[k] = \mathcal{I}, \qquad \mathit{suf}[i] = \Phi\bigl(x_i,\, \mathit{suf}[i{+}1]\bigr),
\]
where $\mathcal{I}$ is $\Phi$'s identity. Then the fold over every child
\emph{except} $c_i$ is exactly
$\mathit{excl}(i) = \Phi\bigl(\mathit{pre}[i],\, \mathit{suf}[i{+}1]\bigr)$.
\end{definition}

\begin{warningbox}[The indexing is the whole bug surface]
$\mathit{pre}[i]$ folds children $0 \dots i{-}1$ (exclusive of $i$) and
$\mathit{suf}[i{+}1]$ folds $i{+}1 \dots k{-}1$. ``Except $c_i$'' is
$\Phi(\mathit{pre}[i], \mathit{suf}[i{+}1])$ -- \emph{not}
$\Phi(\mathit{pre}[i], \mathit{suf}[i])$, which double-counts $c_i$, and not
$\Phi(\mathit{pre}[i{+}1], \mathit{suf}[i{+}1])$, which double-counts again.
Write the two base cases $\mathit{pre}[0]=\mathit{suf}[k]=\mathcal{I}$ first.
\end{warningbox}

\begin{theorem}{Exclusive prefix/suffix is $\bigO{\deg(v)}$ per node, $\bigO{n}$ total}{c20-pre-suf-linear}
For every node $v$, \defref{c20-pre-suf} answers all $k=\lvert\mathit{ch}(v)\rvert$
exclusive aggregates after $2k$ combine calls, and running it at every node
costs $\bigO{\sum_v \deg(v)} = \bigO{n}$.
\end{theorem}

\begin{proofbox}[\thmref{c20-pre-suf-linear}: two linear sweeps, then O(1) per child]
$\mathit{pre}[\cdot]$ is filled by one forward sweep ($k$ combines),
$\mathit{suf}[\cdot]$ by one backward sweep ($k$ combines), so $\bigO{k}$ setup.
Each exclusive query is a single combine of two precomputed cells, so $k$
queries cost $\bigO{k}$. Total per node $\bigO{k}=\bigO{\deg(v)}$. Summed over
the tree, $\sum_v \deg(v) = 2(n-1) = \bigO{n}$ (each edge contributes to two
endpoints). Correctness of $\mathit{excl}(i)$: since $\Phi$ is associative,
$\Phi(\mathit{pre}[i],\mathit{suf}[i{+}1]) = \Phi(x_0,\dots,x_{i-1},x_{i+1},
\dots,x_{k-1})$, the fold of all children but $c_i$; $\mathcal{I}$ handles the
empty ends. \qed
\end{proofbox}

\begin{insightbox}[Why not a segment tree?]
An \term{exclusive segment tree} (build once per node over its children, read
each leaf's complement) is the $\bigO{\deg\log\deg}$ generalisation and is
elegant but over-built for a single pass. For a reroot that visits each node
once, the prefix/suffix pair is the same asymptotics with far less code. Use the
segment tree only when the children list is itself dynamic across queries --
that is Chapter~17's (\chref{c28}) territory, not this chapter's.
\end{insightbox}

\subsection{The Algorithm: Reroot with a Prefix/Suffix Interior}

\begin{lstlisting}[caption={Two-DFS rerooting where the sibling aggregate uses prefix/suffix instead of subtraction -- works for any associative Phi, invertible or not.}, label={lst:c20-presuf}]
// Phi may be NON-invertible (max / min / signed product / or).
// down[v] = fold over subtree(v).  up[v] = fold over the complement.
// ans[v]  = Psi(down[v], up[v]).
Val PHI_ID = identity_of_Phi;   // 0 for +, -INF for max, 1 for product, ...
vector<vector<int>> adj;
vector<Val> down_, up_, ans_;

Val phi(Val a, Val b);                       // associative combine (no inverse!)
Val shift(Val aggregate, int v, int c);

void dfs_down(int v, int p) {                         // post-order
    down_[v] = self_value(v);
    for (int c : adj[v]) if (c != p) {
        dfs_down(c, v);
        down_[v] = phi(down_[v], down_[c]);           // fold child
    }
}
void dfs_up(int v, int p) {                           // pre-order
    ans_[v] = psi(down_[v], up_[v]);
    // collect children (all neighbours except parent)
    vector<int> ch;
    for (int c : adj[v]) if (c != p) ch.push_back(c);
    int k = ch.size();
    vector<Val> pre(k + 1, PHI_ID), suf(k + 1, PHI_ID);
    for (int i = 0; i < k; i++)      pre[i + 1] = phi(pre[i], down_[ch[i]]);
    for (int i = k - 1; i >= 0; i--) suf[i]     = phi(down_[ch[i]], suf[i + 1]);
    for (int i = 0; i < k; i++) {
        int c = ch[i];
        // everything not in subtree(c): self, up[v], and all siblings of c
        Val siblings = phi(pre[i], suf[i + 1]);        // EXCLUDE ch[i] -- thm
        Val rest = phi(phi(up_[v], self_value(v)), siblings);
        up_[c] = shift(rest, v, c);
        dfs_up(c, v);
    }
}
// seed: dfs_down(1,0); up_[1] = PHI_ID; dfs_up(1,0);  read ans_[v].
// Total O(sum deg) = O(n): two O(k) sweeps and one O(1) query per child.
\end{lstlisting}

\begin{figure}[H]
\centering
\begin{tikzpicture}[diagram, scale=0.9, transform shape,
    cnode/.style={circle, draw=primaryblue, line width=0.8pt, fill=primaryblue!8,
                  minimum size=7mm, font=\small\bfseries, text=primaryblue!60!black},
    xnode/.style={circle, draw=warningred, line width=0.8pt, fill=warningred!10,
                  minimum size=7mm, font=\small\bfseries, text=warningred!60!black}]
  \node[xnode] (v)  at (0,0)  {$v$};
  \node[cnode] (c0) at (-3.2,-1.6) {$c_0$};
  \node[cnode] (c1) at (-1.1,-1.6) {$c_1$};
  \node[cnode] (c2) at ( 1.1,-1.6) {$c_2$};
  \node[cnode] (c3) at ( 3.2,-1.6) {$c_3$};
  \draw[tedge] (v)--(c0) (v)--(c1) (v)--(c2) (v)--(c3);
  % prefix arrows on top
  \node[annot, font=\scriptsize, text=successgreen!50!black, align=left] at (-3.9,0.7)
    {$\mathit{pre}[0]{=}\mathcal{I}$};
  \draw[->, successgreen!50!black, line width=0.7pt] (c0.north) -- (c1.north) -- (c2.north) -- (c3.north);
  \node[annot, font=\scriptsize, text=successgreen!50!black] at (0,0.7)
    {$\mathit{pre}[i{+}1]=\Phi(\mathit{pre}[i],x_i)$};
  % suffix arrow on the row of children
  \draw[<-, difforange!70!black, line width=0.7pt] (c0.south) -- (c1.south) -- (c2.south) -- (c3.south);
  \node[annot, font=\scriptsize, text=difforange!70!black] at (0,-2.5)
    {$\mathit{suf}[i]=\Phi(x_i,\mathit{suf}[i{+}1])$,  $\mathit{suf}[k]{=}\mathcal{I}$};
  % exclusion example: everything but c2
  \draw[dashed, warningred, line width=0.8pt, rounded corners]
    (-1.75,-1.15) rectangle (3.85,-2.05);
  \node[annot, font=\scriptsize, text=warningred, anchor=west] at (4.0,-1.6)
    {$\mathit{excl}(2)=\Phi(\mathit{pre}[2],\mathit{suf}[3])$ excl.\ $c_2$};
\end{tikzpicture}
\caption{The exclusive prefix/suffix aggregate at a node $v$ with children
$c_0..c_3$. The top arrow is the forward prefix fold, the bottom the backward
suffix fold; combining $\mathit{pre}[i]$ (children before $i$) with
$\mathit{suf}[i{+}1]$ (children after $i$) yields the fold of \emph{every} child
but $c_i$ in $\bigO{1}$ -- the replacement for ``subtract the child'' when
$\Phi$ has no inverse.}\label{fig:c20-presuf}
\end{figure}

\subsection{Worked Example: a Signed Sum That Flips on Every Edge}

\begin{example}{CF 960E's signed reroot on the running tree $T$}{c20-presuf-run}
Reuse $T$ (root $1$, $n=10$) with node values $V=[2,1,-1,3,-2,1,0,-1,2,1]$.
CF~960E sums, over all $n^2$ directed paths, the alternating-sign node-value
sum of the path. The per-node machinery is: $d[v]=\sum_{u\in\subtree(v)}
(-1)^{\dist(v,u)}$ (the down fold, $d[v]=1-\sum_c d[c]$), $\mathit{tot}[v]=
\sum_u (-1)^{\dist(v,u)}$ (which obeys the reroot flip $\mathit{tot}[c]=
-\mathit{tot}[v]$, since moving the root across one edge changes every distance
by an odd number), and $S[v]=\sum_c d[c]\size(c)$. Then the node's coefficient is
$Q(v) = n\,\mathit{tot}[v] - \bigl(-S[v] + (\mathit{tot}[v]-d[v])(n-\size(v))\bigr)$
and the answer is $\sum_v V[v]\cdot Q(v)$. All values are hand-computed below.
\end{example}

\begin{center}
\footnotesize
\renewcommand{\arraystretch}{1.15}
\begin{tabular}{l|cccccccccc}
\toprule
node $v$ & 1 & 2 & 3 & 4 & 5 & 6 & 7 & 8 & 9 & 10 \\
\midrule
$\size(v)$ & 10 & 4 & 5 & 1 & 2 & 1 & 3 & 1 & 2 & 1 \\
$V[v]$ & 2 & 1 & -1 & 3 & -2 & 1 & 0 & -1 & 2 & 1 \\
$d[v]$ & 2 & 0 & -1 & 1 & 0 & 1 & 1 & 1 & 0 & 1 \\
$\mathit{tot}[v]$ & 2 & -2 & -2 & 2 & 2 & 2 & 2 & -2 & -2 & 2 \\
$S[v]$ & -5 & 1 & 4 & 0 & 1 & 0 & 0 & 0 & 1 & 0 \\
$Q(v)$ & 15 & -7 & -11 & 11 & 5 & 11 & 13 & 7 & -3 & 11 \\
$V\!\cdot\!Q$ & 30 & -7 & 11 & 33 & -10 & 11 & 0 & -7 & -6 & 11 \\
\bottomrule
\end{tabular}
\end{center}

\noindent Sanity: $d[1]=2$, so $\mathit{tot}[1]=2$ and $\mathit{tot}$ then flips
sign across each edge. For node~$1$: $P(1)=-(-5)+0=5$, hence
$Q(1)=10{\cdot}2-5=15$. For node~$2$: $P(2)=-1+(-2){\cdot}6=-13$, hence
$Q(2)=10{\cdot}(-2)-(-13)=-7$. The final answer is $\sum_v V[v]\,Q(v)=66$.

\begin{tracebox}[Computing $d$ (down) and $Q$ for two nodes of $T$]
Down fold $d[v]=1-\sum_{c}d[c]$ (post-order):
leaves $d[4]{=}d[6]{=}d[8]{=}d[10]{=}1$; $d[9]=1-d[10]=0$; $d[5]=1-d[8]=0$;
$d[7]=1-d[9]=1$; $d[2]=1-(d[4]{+}d[5])=1-(1{+}0)=0$;
$d[3]=1-(d[6]{+}d[7])=1-(1{+}1)=-1$; $d[1]=1-(d[2]{+}d[3])=1-(0{-}1)=2$. \checkmark\\
Node $2$: $S[2]=d[4]\size(4)+d[5]\size(5)=1{\cdot}1+0{\cdot}2=1$;
$\mathit{tot}[2]=-\mathit{tot}[1]=-2$; $U[2]=\mathit{tot}[2]-d[2]=-2$;
$P[2]=-S[2]+U[2](n-\size 2)=-1+(-2){\cdot}6=-13$;
$Q(2)=10{\cdot}(-2)-(-13)=-20+13=-7$. \checkmark\\
Node $1$: $S[1]=d[2]\size 2+d[3]\size 3=0{\cdot}4+(-1){\cdot}5=-5$;
$U[1]=0$; $P[1]=-(-5)+0=5$; $Q(1)=10{\cdot}2-5=15$. \checkmark
\end{tracebox}

\subsection{Variations and Twists}

\begin{itemize}[nosep]
  \item \textbf{Pure $\max$ reroot:} $\mathit{excl}$ gives the best sibling so
        $v$'s child can inherit ``best excluding self'' -- CF~1324F /
        CF~342E style up-values. \hfill\lv{4}
  \item \textbf{Product / sign combine:} identity $1$, and a per-edge factor
        (this problem's flip) -- keep $\mathit{pre}/\mathit{suf}$ in
        \lstinline{long long} mod $10^9{+}7$. \hfill\lv{4}--\lv{5}
  \item \textbf{Invertible special case:} when $\Phi$ is $+$ or $\oplus$,
        prefix/suffix degenerates to $\mathit{down}[v]\ominus\mathit{down}[c]$
        (\secref{c20-updown}) -- faster but less general. \hfill\lv{3}
  \item \textbf{Exclusive segment tree:} the same idea at
        $\bigO{\deg\log\deg}$, needed only for dynamic child sets --
        forward-cited to Chapter~17 (\chref{c28}). \hfill\lv{5}
\end{itemize}

\subsection{Common Pitfalls}

\begin{warningbox}[Trap 1 -- double-counted ``except $c$'' from wrong indices]
Using $\mathit{suf}[i]$ (or $\mathit{pre}[i{+}1]$) in the exclusion re-adds
$x_i$; the error hides whenever a node has $2$ children and bites on wide nodes.
Memorise: the neighbour cells of child $i$ are $\mathit{pre}[i]$ and
$\mathit{suf}[i{+}1]$.
\end{warningbox}
\begin{warningbox}[Trap 2 -- wrong identity for a $\max$-fold]
$\mathit{pre}[0]$ and $\mathit{suf}[k]$ must be $-\infty$ for $\max$, $0$ for
$+$, $1$ for $\times$, $\mathit{false}$ for $\mathit{or}$. Seed $\mathit{pre}[0]=0$
for a $\max$ over negative values and the empty prefix silently wins.
\end{warningbox}
\begin{warningbox}[Trap 3 -- mod arithmetic on signed sums]
$Q(v)$ and the $\sum V\!\cdot\!Q$ accumulate negatives; reduce as
\lstinline{((x % MOD) + MOD) % MOD}. Also $d$ and $\mathit{tot}$ are signed
counts, not sizes -- never take their absolute value.
\end{warningbox}

\subsection{Problem Lab}

\solvedproblem{c20-presuf-solved}{Alternating Tree}{Codeforces 960E}
  {https://codeforces.com/problemset/problem/960/E}{CF 2300 · L5}
  {signed-sum reroot, per-edge sign flip, exclusive prefix/suffix sibling
   aggregate, modulo counting, long long}

\begin{problemstatement}[Alternating Tree (statement and sample verified against the live problem)]
A tree of $n$ nodes ($2 \le n \le 2\cdot10^{5}$), node $i$ has value $V_i$
($-10^9 \le V_i \le 10^9$). For a simple path $u_1 \to u_m$ (nodes $u_1,\dots,u_m$,
possibly $u_1 = u_m$) the alternating function is
$A(u_1,u_m) = \sum_{i=1}^{m} (-1)^{i+1} V_{u_i}$. Compute the sum of
$A(u_1,u_m)$ over \emph{all directed} paths -- $(u,v)$ and $(v,u)$ are
different, and single-node paths count -- modulo $10^9+7$.
\textbf{Sample}: $n=4$, $V=[-4,1,5,-2]$, star edges $1\text{--}2,1\text{--}3,
1\text{--}4$; the note lists all $16$ directed paths and their sum is
\texttt{40}.
\end{problemstatement}

\begin{attemptgate}[Attempt it yourself first -- \lv{5}, give yourself 40 minutes]
Reorder the double sum. (a)~A node $w$ sitting at position $i$ on the path
$u\!\to\!v$ contributes $(-1)^{i+1}V_w$, and $i = \dist(u,w)+1$, so its sign is
$(-1)^{\dist(u,w)}$. (b)~So the whole answer is $\sum_w V_w \cdot Q(w)$ where
$Q(w)$ is a purely structural count. Work out what $Q(w)$ counts.
\peek{For a fixed $w$, over which ordered pairs $(u,v)$ does $w$ lie on the
      path, and what is $w$'s sign there?}
\nudge{The sign of $w$ in $A(u,v)$ is $(-1)^{\dist(u,w)}$. $w$ is on the path
      iff $u,v$ are in \emph{different} branches of $T - w$ or one equals $w$.
      Split $Q(w)$ into $u$-contributions per branch $b$ with $D_b=\sum_{v\in b}
      (-1)^{\dist(w,v)}$.}
\fullhint{$Q(w)$ is the signed count (ordered pairs through $w$ with $w$ at
shorter start-distance even) $-$ (odd). For branch $b$, pairs choosing $u\in b$
and $v$ outside $b$ number $\size_b\,(n-\size_b)$, splitting by parity into
$D_b(n-\size_b)$. Adding the $u=w$ case ($n$ pairs) and simplifying gives
$Q(w) = n + \sum_b D_b(n-\size_b) = n\,\mathit{tot}[w] - \sum_b D_b\,\size_b$.
Now $D_b$ for a child branch is $-d[c]$, and for the parent branch
$\mathit{tot}[w]-d[w]$, where $d[v]=1-\sum_c d[c]$ and $\mathit{tot}[c]=
-\mathit{tot}[v]$ (reroot sign flip). Compute $d,\mathit{tot},\size$ by two DFS
passes (prefix/suffix handles the ``all branches but this one'' combine
$\sum_b$ generically), then return $\sum_w V_w Q(w)$ reduced mod $10^9{+}7$.}
\end{attemptgate}

\solstep{Understanding the Problem}
The definition sums an alternating quantity over $n^2$ directed paths -- far too
many to enumerate. The rescue is linearity: swap the order of summation so we
walk \emph{nodes}, not paths. Each node $w$ with value $V_w$ appears on many
paths, and on a given $u\!\to\!v$ path it sits at position $i=\dist(u,w)+1$, so
it contributes $(-1)^{\dist(u,w)}V_w$. Factoring $V_w$ out, everything else is
structural: $Q(w) = \sum_{u,v \text{ through } w} (-1)^{\dist(u,w)}$. The task
becomes computing one signed count per node -- the every-node quantity the
two-DFS template exists for. The combine that forms each branch aggregate
involves a sign that flips across every edge, so even though the arithmetic is
addition, the ``everything but this branch'' pieces enter a product
$D_b(n-\size_b)$, and we compute them with \defref{c20-pre-suf}'s exclusive
aggregate for generality.

\solstep{Key Observations}
\begin{observation}
$\sum_{\text{directed paths}} A = \sum_w V_w\,Q(w)$ with
$Q(w)=\sum_{(u,v):\,w\in\text{path}} (-1)^{\dist(u,w)}$ -- reordering the double
sum; the sign of $w$ depends only on $\dist(u,w)$'s parity.
\end{observation}
\begin{observation}
$\mathit{tot}[v]=\sum_{u\in V}(-1)^{\dist(v,u)}$ obeys $\mathit{tot}[c]=
-\mathit{tot}[v]$ for adjacent $v,c$: sliding the root one edge changes the
distance to \emph{every} node by an odd number, flipping every term's sign. This
is the per-edge sign flip of the reroot.
\end{observation}
\begin{observation}
$Q(w)=n\,\mathit{tot}[w]-\sum_b D_b\size_b$, where the sum is over branches $b$
of $T-w$, $D_b=\sum_{v\in b}(-1)^{\dist(w,v)}$ and $\size_b$ is the branch size
-- the $u=w$ pairs contribute $n$, the rest regroup into the branch products.
Child branch $D_b=-d[c]$, parent branch $D_b=\mathit{tot}[w]-d[w]$.
\end{observation}

\solstep{Think It Through With Me}
\begin{thinkbox}
\deadend{Enumerate all $n^2$ ordered pairs and walk each path with binary
        lifting to evaluate its alternating sum.}
{$\bigO{n^2 \log n}$, and $n=2\cdot10^5$ makes that $\sim10^{11}$ steps --
certain \tle. It also ignores the sign-parity structure entirely.}
\deadend{Treat $d[v]=1-\sum_c d[c]$ as a plain component size and skip the
        reroot, reading only $d[1]$.}
{$Q(w)$ needs the signed count $\mathit{tot}[w]$ and the branch aggregate at
\emph{every} node, not just the reference root. Reading one node answers nothing
about the others -- you must reroot, and the reroot's edge move is a sign flip
because parity is what changes.}
\aha{Linearise into $\sum_w V_w Q(w)$. Reroot $\mathit{tot}$ with the flip
$\mathit{tot}[c]=-\mathit{tot}[v]$; get $D_b,\size_b$ per branch through a
prefix/suffix exclusive combine; $Q(w)=n\,\mathit{tot}[w]-\sum_b D_b\size_b$.
That is \coderef{c20-presuf} with $\Phi=+$ carrying a per-edge sign.}
\end{thinkbox}

\solstep{Solution}
Root at $1$. Post-order: $\size[v]=1+\sum_c\size[c]$, $d[v]=1-\sum_c d[c]$,
$S[v]=\sum_c d[c]\size[c]$. $\mathit{tot}[1]=d[1]$. Pre-order: $\mathit{tot}[c]=
-\mathit{tot}[v]$. Then $U[w]=\mathit{tot}[w]-d[w]$,
$P[w]=-S[w]+U[w](n-\size[w])$, $Q[w]=n\,\mathit{tot}[w]-P[w]$. Answer
$=\bigl(\sum_w V_w\,Q_w\bigr)\bmod(10^9+7)$. (The $\sum_b D_b\size_b$ is the
exclusive over-branches aggregate \defref{c20-pre-suf} evaluates generically;
here it collapses to $P[w]$ since all but one branch are children.)

\solstep{Proof of Correctness}
\begin{proofbox}[\coderef{c20-960e-soln} returns the directed-path alternating sum]
\emph{Reordering.} By linearity $\sum_{u,v}A(u,v)=\sum_w V_w Q(w)$ with $Q(w)$ as
in Observation~1 (each occurrence of $w$ carries exactly $(-1)^{\dist(u,w)}$).
\emph{$\mathit{tot}$ flip.} For adjacent $v,c$ and any $u$, the paths $v\!\to\!u$
and $c\!\to\!u$ differ by $\pm1$ edge, so $\dist(c,u)\equiv \dist(v,u)+1\pmod 2$;
hence $(-1)^{\dist(c,u)}=-(-1)^{\dist(v,u)}$ for every $u$, and
$\mathit{tot}[c]=-\mathit{tot}[v]$. \emph{$Q$ formula.} Split $Q(w)$ by whether
$u=w$ (contributes $+1$ times the $n$ choices of $v$: $n$) or $u$ lies in branch
$b$ (then $v$ may be any node outside $b$ -- count $n-\size_b$ -- and the sign is
$(-1)^{\dist(w,u)}$): $Q(w)=n+\sum_b D_b(n-\size_b)$. Since $\sum_b D_b=
\mathit{tot}[w]-1$ and $\sum_b(n-\size_b)$ groups, this rearranges to
$n\,\mathit{tot}[w]-\sum_b D_b\size_b$. Child $D_b=-d[c]$ (signs measured from
$w$ one step further than from $c$), parent $D_b=\mathit{tot}[w]-d[w]$; the
branch-size product $\sum_b D_b\size_b$ is exactly the exclusive aggregate of
\defref{c20-pre-suf}, computed here as $-S[w]+U[w](n-\size[w])$. Combining,
$Q[w]=n\,\mathit{tot}[w]-P[w]$. Summing $V_wQ_w$ and reducing mod $10^9+7$ gives
the required total. \qed
\end{proofbox}

\solstep{Implementation}
\begin{lstlisting}[caption={Codeforces 960E -- signed-sum reroot (tot flip per edge) with the branch product aggregate; O(n), mod 1e9+7}, label={lst:c20-960e-soln}]
#include <bits/stdc++.h>
using namespace std;
typedef long long ll;
const ll MOD = 1000000007LL;
int n;
vector<ll> val, sub, d, S, tot, Q;
vector<vector<int>> adj;

void dfs_down(int v, int p) {            // post-order: size, signed d, S
    sub[v] = 1; d[v] = 1; S[v] = 0;
    for (int c : adj[v]) if (c != p) {
        dfs_down(c, v);
        sub[v] += sub[c];
        d[v]   -= d[c];                  // d[v] = 1 - sum_child d[c]
        S[v]   += d[c] * sub[c];         // sum over child branches
    }
}
void dfs_up(int v, int p) {              // pre-order: tot flip across edges
    for (int c : adj[v]) if (c != p) {
        tot[c] = -tot[v];                // per-edge sign flip, Obs 2
        dfs_up(c, v);
    }
}
ll norm(ll x){ x %= MOD; if (x < 0) x += MOD; return x; }
int main() {
    scanf("%d", &n);
    adj.assign(n + 1, {});
    val.assign(n + 1, 0);
    for (int i = 1; i <= n; i++) { long long x; scanf("%lld", &x); val[i] = x; }
    for (int i = 0; i < n - 1; i++) {
        int u, v; scanf("%d %d", &u, &v);
        adj[u].push_back(v); adj[v].push_back(u);
    }
    sub.resize(n+1); d.resize(n+1); S.resize(n+1); tot.resize(n+1); Q.resize(n+1);
    dfs_down(1, 0);
    tot[1] = d[1];
    dfs_up(1, 0);
    ll ans = 0;
    for (int v = 1; v <= n; v++) {
        ll U   = tot[v] - d[v];                        // parent-branch D
        ll P   = -S[v] + U * (ll)(n - sub[v]);         // sum_b D_b * size_b
        Q[v]   = (ll)n * tot[v] - P;                   // coefficient of val[v]
        ans   += norm(val[v]) * norm(Q[v]);
    }
    printf("%lld\n", norm(ans));
    return 0;
}
\end{lstlisting}

\solstep{Trace}
\begin{tracebox}[Official sample: $n=4$, values $(-4,1,5,-2)$, star centred at $1$ $\Rightarrow$ 40]
Leaves $2,3,4$: $d=1,\size=1,S=0$. $d[1]=1-(1{+}1{+}1)=-2$, $\size[1]=4$,
$S[1]=1{+}1{+}1=3$. $\mathit{tot}[1]=d[1]=-2$; $\mathit{tot}[2]=\mathit{tot}[3]=
\mathit{tot}[4]=-\mathit{tot}[1]=2$.\\
$v=1$: $U=0$; $P=-3+0\cdot0=-3$; $Q=4(-2)-(-3)=-5$; term $(-4)(-5)=20$.\\
$v=2$: $U=2-1=1$; $P=0+1\cdot3=3$; $Q=4(2)-3=5$; term $1\cdot5=5$.\\
$v=3$: $Q=5$; term $5\cdot5=25$. \quad $v=4$: $Q=5$; term $(-2)(5)=-10$.\\
Sum $=20+5+25-10=\boxed{40}$ \checkmark\ matches CF's sample (note lists all
$16$ directed paths summing to $40$).
\end{tracebox}

\solstep{Complexity}
\begin{complexitybox}
Time $\bigO{n}$: one post-order and one pre-order pass; every node's branch
aggregate is $O(\deg)$ with the prefix/suffix idea (or $\bigO{1}$ after the
$S[w]$ / $U[w]$ regrouping used in \coderef{c20-960e-soln}). Memory $\bigO{n}$.
The accumulator is mod $10^9{+}7$ with \lstinline{long long} intermediates
($\mathit{tot}$ is signed, so \lstinline{val*Q} is up to $10^{9}\cdot10^{14}$
before reduction -- take \lstinline{norm} on both factors, as shown). Stack
depth $h$; iterate with \coderef{c09-iter} at $n\le2\cdot10^{5}$. Verdict: \ac;
the archetypal \emph{non-invertible combine} reroot -- a per-edge sign flip that
makes subtraction illegal and pays for the prefix/suffix machinery.
\end{complexitybox}

\solstep{Pitfalls and Variants}
\begin{warningbox}
Ways to lose CF~960E: (i)~reducing $d,\mathit{tot}$ mod $10^9{+}7$ \emph{as
unsigned} counts -- they are signed; use \lstinline{norm}. (ii)~forgetting the
$u=v$ single-node path (contributes $V_u$ with sign $+1$, already inside $Q$ via
the $n\,\mathit{tot}$ term -- verify a $1$-node sanity). (iii)~treating
$\mathit{tot}$ as a size and skipping the flip. (iv)~overflow before the
\lstinline{norm}: reduce both factors. Variants: a $\max$-valued every-root
statistic (CF~1324F / CF~342E) uses \coderef{c20-presuf} verbatim with
$\Phi=\max$ and $\mathcal{I}=-\infty$; a ``product of subtree values for every
root'' reroot needs $\mathcal{I}=1$ and modular inverses only where the combine
permits them.
\end{warningbox}

\subsection{Practice Problems}

\begin{stuckbox}
\begin{itemize}[nosep, leftmargin=1.4em]
  \item Did you reorder the double sum into $\sum_w V_w Q(w)$ before trying to
        walk paths?
  \item Do you have the reroot identity $\mathit{tot}[c]=-\mathit{tot}[v]$ (sign
        flip) rather than a size transfer?
  \item Is the ``all branches but $b$'' aggregate built as
        $\Phi(\mathit{pre}[i],\mathit{suf}[i{+}1])$, and are the signed counts
        reduced with a \lstinline{norm}?
\end{itemize}
\end{stuckbox}

\begin{practiceset}
  \practice{Casinos and Travel}{Codeforces 852E}
    {https://codeforces.com/problemset/problem/852/E}{CF 2100 · L5}
    {John starts at a chosen city and each day moves to an unvisited neighbour,
    stopping when stuck; a casino toggles his mood (good $\leftrightarrow$ bad),
    and he must end \emph{good} on \emph{every} possible walk. Count, mod
    $10^9{+}7$, the ways to pick the start city and a casino subset. A parity
    DP over the rooted tree -- $\mathit{down}[v]$ and a matching
    ``$\mathit{up}$/all-branches-but-this'' combine decide whether each node can
    be the start, so the ``for every node'' reading uses
    \defref{c20-pre-suf}'s exclusive aggregate. Samples (verified): $n=2$ path
    $\to 4$; $n=3$ path $\to 10$.
    }
\end{practiceset}

\paragraph{Back-cites (URLs live elsewhere in the book).}
\begin{itemize}[nosep,leftmargin=1.4em]
  \item \textbf{CF~1324F ``Maximum White Subtree''} (\chref{c10} owns the URL):
        the $\max$-reroot that needs the exclusive sibling fold without an
        inverse -- do it with \coderef{c20-presuf} and $\Phi=\max$.
  \item \textbf{CF~342E ``Cueroten''} (\chref{c09} owns the URL): the
        distance/aggregate-per-root family, a direct prefix/suffix consumer.
\end{itemize}

\begin{recallbox}[10-second refresher -- 15.3 Prefix/suffix merge]
When $\Phi$ has no inverse, get ``all children but $c_i$'' from
$\mathit{excl}(i)=\Phi(\mathit{pre}[i],\mathit{suf}[i{+}1])$ after two
$\bigO{\deg}$ sweeps -- $\bigO{n}$ total (\thmref{c20-pre-suf-linear}); the
identity $\mathcal{I}$ depends on $\Phi$. CF~960E (\probref{c20-presuf-solved})
reroots a signed sum: reorder to $\sum_w V_wQ(w)$, use the per-edge flip
$\mathit{tot}[c]=-\mathit{tot}[v]$, and form the branch products with the
exclusive combine. On $T$ (values $V$ above) the answer is $66$; the official
star gives $40$.
\selftest{
  \item For $\Phi=\max$ over negative values, what are $\mathit{pre}[0]$ and
        $\mathit{suf}[k]$?
  \item Why does moving the root one edge flip the sign of $\mathit{tot}$, not
        its magnitude?
  \item Give the off-by-one that double-counts $c_i$ in the exclusion.
  \item What makes a plain subtraction illegal for CF~960E's combine?
  \item State the $u=w$ (single-node path) contribution to $Q(w)$.
}
\end{recallbox}
